\documentclass[11pt]{article}
\usepackage[utf8]{inputenc}
\usepackage[T1]{fontenc}
\usepackage{amsmath,amssymb,mathtools}
\usepackage[margin=1in]{geometry}
\usepackage{hyperref}
\usepackage{array}
\usepackage{parskip}
\setlength{\parindent}{0pt}
\hypersetup{colorlinks=true,linkcolor=blue,urlcolor=blue}
\title{Lecture 5 — Basic System Properties}
\author{}
\date{}
\begin{document}
\maketitle

\textbf{Course:} CEC 315 — Signals and Systems (Spring 2026)
\textbf{Instructor:} Jianhua Liu
\textbf{Source PDF:} \texttt{all\_lectures/cec315-lctr05-basic-sys-properties.pdf}
\textbf{Exam coverage:} Exam 1

\textbf{Lctr 5: Basic System Properties}

Jianhua Liu

Spring 2026

\textbf{Focus:} Section 1.6 (Pages 44 to 55)

\section*{5.1 Memory vs. Memoryless Systems}

\begin{quote}
\textbf{Why This Matters}

Every system you will encounter in this course — and in your engineering career — can be classified by a small set of "basic properties": memory, invertibility, causality, stability, time invariance, and linearity. These properties aren't just vocabulary; they determine which mathematical tools you can use to analyze the system. Only \textbf{linear time-invariant (LTI)} systems can be analyzed with convolution, Fourier, Laplace, and $z$-transforms. So the first step in analyzing any system is to check these properties.
\end{quote}

\subsection*{5.1.1 Memoryless Systems}

\begin{quote}
\textbf{Definition — Memoryless System}

A system is \textbf{memoryless} if the output at any time depends \textbf{only on the input at that same time}.

Mathematically: $y(t)$ depends only on $x(t)$ (CT), or $y[n]$ depends only on $x[n]$ (DT).
\end{quote}

\textbf{Example:} The \textbf{ReLU} (Rectified Linear Unit) activator used in neural networks:

\[y[n] = \begin{cases} x[n], & x[n] \geq 0 \\ 0, & \text{otherwise}\end{cases}\]

For any given $n$, the output $y[n]$ is computed purely from $x[n]$ — no past or future samples are involved. ReLU is memoryless.

\textbf{Other memoryless examples:}

\begin{itemize}
  \item $y(t) = 5 x(t)$ (pure scaling)
  \item $y(t) = x^2(t)$ (squaring)
  \item $y[n] = \cos(x[n])$
\end{itemize}

\subsection*{5.1.2 Systems With Memory}

\begin{quote}
\textbf{Definition — System With Memory}

A system has \textbf{memory} if the output depends on past or future values of the input (or on past values of the output itself).
\end{quote}

\textbf{Examples:}

\begin{itemize}
  \item \textbf{Recursive (IIR) system:} $y_1[n] = x[n] + \alpha\,y_1[n-1]$ The current output depends on the previous output $y_1[n-1]$ in addition to the current input.
  \item \textbf{Moving-average / FIR system:} $y_2[n] = x[n] - \beta\,x[n-1]$ The current output depends on the previous input $x[n-1]$ in addition to the current input.
  \item \textbf{Capacitor voltage:} $v_C(t) = \frac{1}{C}\int_{-\infty}^{t} i(\tau)\,d\tau$ The voltage at time $t$ remembers all past current.
\end{itemize}

\begin{quote}
\textbf{Key Insight}

Any system that contains an integrator, a delay element, a difference, or a recursion has memory. Memoryless systems are algebraic (or at most, instantaneous functions of the current sample).
\end{quote}

\section*{5.2 Invertibility and Inverse Systems}

\subsection*{5.2.1 Definitions of Inverse Systems}

\begin{quote}
\textbf{Definition — Invertible System}

Two equivalent definitions:

1. A system is \textbf{invertible} if distinct inputs lead to distinct outputs (the input-output map is one-to-one).
2. A system $T_1(\cdot)$ is invertible if there exists another system $T_2(\cdot)$ such that the impulse response of the cascade $T_2\!\circ\!T_1$ equals the unit impulse. In this case, $T_1$ and $T_2$ are inverse systems of each other.
\end{quote}

\textit{[Figure: Block diagram — $x[n] \to T_1 \to w[n] \to T_2 \to x[n]$. The overall cascade is the identity system; its impulse response is $\delta[n]$.]}

\textbf{Remark.} In many real applications we cannot find an exact inverse, but we can find a system such that the cascade is an \textbf{all-pass filter} (magnitude response equal to 1 for all frequencies). This technique underlies \textbf{minimum-phase filter} design and is explored in CEC 410.

\subsection*{5.2.2 Finding an Inverse System (Preview via $z$-Transform)}

Consider the two DT systems introduced above:

\[\text{System 1: } y_1[n] = x[n] + \alpha\,y_1[n-1]\]
\[\text{System 2: } y_2[n] = x[n] - \beta\,x[n-1]\]

Using the $z$-transform (to be formally introduced later in the course), their system functions are:

\[H_1(z) = \frac{1}{1 - \alpha z^{-1}}, \qquad H_2(z) = 1 + \beta z^{-1}.\]

Setting $\alpha = -\beta$ gives:

\[H_1(z)\,H_2(z) = \frac{1}{1 + \beta z^{-1}}\,(1 + \beta z^{-1}) = 1.\]

Since the $z$-transform of $\delta[n]$ is $1$, the two systems are indeed inverses of each other. This preview illustrates \textit{why} transform-domain analysis is powerful: inversion becomes algebraic.

\section*{5.3 Causality}

\subsection*{5.3.1 Definition of a Causal System}

\begin{quote}
\textbf{Definition — Causal System}

Two equivalent definitions:

1. A system is \textbf{causal} if the output at any time depends only on \textbf{present and past} values of the input.
2. A system is \textbf{causal} if the output at any time does \textbf{not depend on future input}.
\end{quote}

\subsection*{5.3.2 Quick Test via the Impulse Response}

For an LTI system, there is an easy test:

\[\boxed{\,\text{LTI system is causal} \iff h[n] = 0 \text{ for } n < 0 \text{ (or } h(t) = 0 \text{ for } t < 0\text{)}\,}\]

\textbf{Why?} The impulse response is the system's output when the input is $\delta[n]$. The impulse "happens" at $n=0$; it has no energy at $n<0$. A causal system cannot produce a nonzero response before the stimulus arrives.

\subsection*{5.3.3 Examples of Causal / Non-Causal Systems}

\begin{center}
\begin{tabular}{|l|l|l|}
\hline
\textbf{System} & \textbf{Causal?} & \textbf{Reasoning} \\
\hline
$y_1[n] = x[n] + \alpha y_1[n-1]$ & \textbf{Yes} & Depends on present input and past output only. \\
\hline
$y_2[n] = x[n] - \beta x[n-1]$ & \textbf{Yes} & Depends on present and past input only. \\
\hline
$y_3[n] = x[n] - \beta x[n+1]$ & \textbf{No} & Depends on future input $x[n+1]$. \\
\hline
ReLU $y[n] = \max(x[n],0)$ & \textbf{Yes} & Only depends on current input. \\
\hline
\end{tabular}
\end{center}

\begin{quote}
\textbf{Key Insight}

All \textbf{real-time} physical systems are causal (time flows forward). Non-causal systems arise naturally in \textit{offline} processing: audio editing, image processing, batch-mode speech recognition, etc. Any system that "peeks" at future samples is non-causal.
\end{quote}

\section*{5.4 Stability}

\subsection*{5.4.1 Definition of a Stable System}

Informally, a stable system is one where \textbf{small inputs do not lead to diverging responses}.

\begin{quote}
\textbf{Definition — BIBO Stability (Bounded-Input, Bounded-Output)}

A system is \textbf{BIBO stable} if every bounded input produces a bounded output:

\[|x(t)| \leq B_x < \infty \;\;\Longrightarrow\;\; |y(t)| \leq B_y < \infty\]

for some constants $B_x, B_y$.
\end{quote}

\subsection*{5.4.2 Examples of Stable / Unstable Systems}

\begin{itemize}
  \item \textbf{Pendulum} (hanging down): stable — small pushes produce bounded swings.
  \item \textbf{Inverted pendulum} (balanced on its tip): unstable — any perturbation grows without bound.
  \item $y[n] = 0.5\,y[n-1] + x[n]$: stable (pole inside unit circle, as we will see later).
  \item $y[n] = 2\,y[n-1] + x[n]$: unstable (pole outside unit circle).
  \item Integrator $y(t) = \int_{-\infty}^t x(\tau)\,d\tau$: \textbf{not} BIBO stable. Feeding in the bounded input $x(t) = u(t)$ (the unit step) gives $y(t) = t\,u(t)$, which is unbounded.
\end{itemize}

\begin{quote}
\textbf{Key Insight}

Stability is the single most important property for control and filter design. An unstable system is useless in practice — it will saturate, oscillate, or destroy hardware. We will later derive stability conditions in terms of poles (Laplace: LHP; $z$-transform: inside unit circle) and in terms of $h(t)$ or $h[n]$ (absolutely integrable / summable).
\end{quote}

\section*{5.5 Time Invariance}

\subsection*{5.5.1 Definition of Time Invariance}

\begin{quote}
\textbf{Definition — Time-Invariant System}

A system is \textbf{time-invariant} if its behavior and characteristics do \textbf{not change over time}. Equivalently, a \textbf{time shift in the input} produces an \textbf{identical time shift in the output}.

Formally: if $x(t) \to y(t)$, then $x(t - t_0) \to y(t - t_0)$ for every $t_0$.
\end{quote}

\subsection*{5.5.2 Verifying Time Invariance}

\textbf{Procedure:}

\begin{enumerate}
  \item Let $y_1(t)$ be the output when the input is $x_1(t)$.
  \item Shift the input: $x_2(t) = x_1(t - t_0)$. Compute $y_2(t)$ directly from the system equation.
  \item Compute $y_1(t - t_0)$ — that is, shift the \textit{output} of step 1 by $t_0$.
  \item If $y_2(t) = y_1(t - t_0)$ for every $t_0$ and every input, the system is time invariant.
\end{enumerate}

\textbf{Example (time invariant):} $y(t) = 3\,x(t) + 2$.
Then $y_2(t) = 3\,x_1(t-t_0) + 2$ and $y_1(t-t_0) = 3\,x_1(t-t_0) + 2$. These match: time invariant.

\textbf{Example (time varying):} $y(t) = t\,x(t)$.
Then $y_2(t) = t\,x_1(t-t_0)$ but $y_1(t-t_0) = (t-t_0)\,x_1(t-t_0)$. These don't match unless $t_0 = 0$: \textbf{time varying}. The explicit $t$ multiplier is a red flag.

\begin{quote}
\textbf{Warning}

Any system equation that contains an \textbf{explicit} dependence on $t$ or $n$ (outside of the input/output signals) is typically time varying. Watch for things like $t x(t)$, $x(2t)$ (time scaling), or coefficients that themselves depend on time.
\end{quote}

\section*{5.6 Linearity}

\subsection*{5.6.1 Definition of Linearity}

\begin{quote}
\textbf{Definition — Linear System}

A system is \textbf{linear} if it possesses the \textbf{superposition} property. Superposition = additivity + homogeneity.
\end{quote}

\subsection*{5.6.2 Additivity}

If $x_1(t) \to y_1(t)$ and $x_2(t) \to y_2(t)$, then:

\[x_1(t) + x_2(t) \;\longrightarrow\; y_1(t) + y_2(t).\]

\subsection*{5.6.3 Scaling (Homogeneity)}

If $x_1(t) \to y_1(t)$, then for every constant $a$ (real or complex):

\[a\,x_1(t) \;\longrightarrow\; a\,y_1(t).\]

\subsection*{5.6.4 Combined Superposition Statement}

\[\boxed{\;a\,x_1(t) + b\,x_2(t) \;\longrightarrow\; a\,y_1(t) + b\,y_2(t)\;}\]

\subsection*{5.6.5 Zero-In / Zero-Out Test (Necessary Condition)}

\begin{quote}
\textbf{Key Insight}

For a linear system, a \textbf{zero input must produce a zero output}. Setting $a = 0$ in homogeneity gives $0 \to 0$. So if you can find any input that's zero but produces a nonzero output (e.g., a constant offset), the system is nonlinear.
\end{quote}

\textbf{Example (nonlinear):} $y(t) = 3\,x(t) + 2$. Zero input gives $y(t) = 2 \neq 0$. \textbf{Nonlinear} (it's \textit{affine}, not linear). Though earlier we found it to be time invariant, linearity fails.

\textbf{Example (linear):} $y(t) = 5\,x(t)$. Obviously linear.

\textbf{Example (nonlinear):} $y(t) = x^2(t)$. Response to $x_1+x_2$ is $(x_1+x_2)^2 = x_1^2 + 2 x_1 x_2 + x_2^2$, not $x_1^2 + x_2^2$. Fails additivity.

\subsection*{5.6.6 ASR (Speech Recognition) Case Study}

Using \textbf{Automatic Speech Recognition (ASR)} as an end-to-end example of how these properties apply to different components of a real system:

\begin{itemize}
  \item \textbf{Memory}
  \item Mel-frequency vector calculation: \textbf{memoryless} (each frame independent).
  \item Phoneme estimation: has memory (context windows, recurrent models).
  \item \textbf{Causality}
  \item A \textbf{streaming} ASR system is causal (must respond in real time).
  \item A \textbf{batch-processing} ASR system is non-causal (can look ahead).
  \item \textbf{Stability}
  \item The overall system is stable (bounded audio produces bounded labels).
  \item \textbf{Linearity}
  \item The entire system is \textbf{nonlinear} (classification is an inherently nonlinear operation).
  \item Individual components can be linear — for example, the \textbf{FFT} inside the mel-frequency extraction is a linear operator.
\end{itemize}

\section*{5.7 Exercises}

\subsection*{5.7.1 Exercise 1: System Classification}

Consider:

\[T_1(x[n]) = \frac{1}{x[n+1]}, \qquad T_2(x[n]) = x[n] + 0.5\,x[n+1].\]

Determine whether each system is: invertible, causal, BIBO stable, time invariant, linear, and memoryless.

\textbf{Analysis of $T_1$:}

\begin{itemize}
  \item \textbf{Memoryless?} No — uses $x[n+1]$.
  \item \textbf{Causal?} No — depends on future input $x[n+1]$.
  \item \textbf{Linear?} No — $T_1(a\,x) = 1/(a\,x[n+1]) = (1/a)\,T_1(x) \neq a\,T_1(x)$. Fails homogeneity.
  \item \textbf{Time invariant?} Yes — shifting the input shifts the output identically.
  \item \textbf{BIBO stable?} No — if $x[n+1] \to 0$ (bounded input going through zero), output diverges.
  \item \textbf{Invertible?} Yes — from $y[n] = 1/x[n+1]$ we recover $x[n+1] = 1/y[n]$, so $x[n] = 1/y[n-1]$.
\end{itemize}

\textbf{Analysis of $T_2$:}

\begin{itemize}
  \item \textbf{Memoryless?} No — uses $x[n+1]$.
  \item \textbf{Causal?} No — depends on $x[n+1]$.
  \item \textbf{Linear?} Yes — sum and scale carry through directly.
  \item \textbf{Time invariant?} Yes.
  \item \textbf{BIBO stable?} Yes — $|y[n]| \leq |x[n]| + 0.5|x[n+1]| \leq 1.5\,B_x$.
  \item \textbf{Invertible?} Yes (it has an inverse, though it requires solving a first-order difference equation).
\end{itemize}

\subsection*{5.7.2 Exercise 2: Signal Sketching (Preparation for Convolution)}

Sketch the following DT signals:

\begin{itemize}
  \item $x_1[n] = (0.5)^n\,u[n]$ — a decaying geometric sequence starting at $n=0$ with value $1$ and halving each step.
  \item $x_2[n] = x_1[n-2]$ — the same shape, shifted right by 2 (starts at $n=2$).
  \item $x_3[n] = x_2[-n]$ — time reversal of $x_2$, giving a sequence that \textit{ends} at $n = -2$ with the same decaying shape mirrored.
\end{itemize}

\textit{[Figure: Three DT stem plots. (a) $x_1[n]$: stems at $n=0,1,2,\ldots$ of heights $1, 0.5, 0.25, 0.125, \ldots$; (b) $x_2[n]$: same heights starting at $n=2$; (c) $x_3[n]$: same heights at $n=-2, -3, -4, \ldots$ (mirrored into negative time).]}

\section*{5.8 Summary}

Key takeaways from this lecture:

\begin{enumerate}
  \item \textbf{Memory:} Output depends only on the current input ($\Rightarrow$ memoryless) vs. past/future input or past output ($\Rightarrow$ memory).
  \item \textbf{Invertibility:} One-to-one map from inputs to outputs; cascade with inverse gives identity $\delta[n]$.
  \item \textbf{Causality:} Output depends only on present and past input. For LTI, equivalent to $h[n] = 0$ for $n<0$.
  \item \textbf{Stability (BIBO):} Bounded input produces bounded output. Unstable systems are unusable in practice.
  \item \textbf{Time invariance:} Shifting the input by $t_0$ shifts the output by exactly $t_0$; equivalently, no explicit $t$ coefficients.
  \item \textbf{Linearity:} Superposition — additivity and scaling. Zero input must give zero output.
\end{enumerate}

The next three lectures will focus on \textbf{LTI systems} (linear \textbf{and} time-invariant) because this is the class for which convolution applies and all the transform machinery of this course is valid.

\section*{5.9 Common Mistakes to Avoid}

\begin{enumerate}
  \item \textbf{Calling affine systems linear.} $y = mx + b$ is \textit{affine}, not linear. Always apply the zero-in / zero-out test.
  \item \textbf{Conflating causality and stability.} They are independent: a system can be causal but unstable, or stable but non-causal.
  \item \textbf{Forgetting to check both additivity AND homogeneity} when testing linearity. Some systems satisfy one but not the other.
  \item \textbf{Assuming time invariance from smooth coefficients.} The telltale sign of time variance is an explicit $t$ or $n$ outside the signal arguments (e.g., $y(t) = t\,x(t)$).
  \item \textbf{Confusing "has memory" with "uses $x[n-1]$".} Both recursive ($y[n-1]$) and delay-line ($x[n-1]$) structures have memory; so do systems with future dependence ($x[n+1]$, though that also makes them non-causal).
\end{enumerate}

Jianhua Liu

\section*{Practice Problems Summary}

\begin{enumerate}
  \item \textbf{Memoryless check:} ReLU $y[n] = \max(x[n],0)$ is memoryless; $y_1[n] = x[n] + \alpha y_1[n-1]$ has memory.
  \item \textbf{Inverse system:} For $H_1(z) = 1/(1-\alpha z^{-1})$ and $H_2(z) = 1 + \beta z^{-1}$, choosing $\alpha = -\beta$ gives $H_1 H_2 = 1$, confirming they are inverses.
  \item \textbf{Causality check:} $y_3[n] = x[n] - \beta x[n+1]$ is non-causal (uses future input); all other earlier examples are causal.
  \item \textbf{Stability intuition:} Pendulum stable; inverted pendulum unstable. Integrator fed by $u(t)$ gives $t\,u(t)$, proving it is not BIBO stable.
  \item \textbf{Time invariance check:} $y(t) = t\,x(t)$ is time varying due to the explicit $t$ coefficient.
  \item \textbf{Linearity check:} $y(t) = 3 x(t) + 2$ is nonlinear (nonzero output for zero input).
  \item \textbf{Exercise 1:} $T_1(x[n]) = 1/x[n+1]$ is non-causal, nonlinear, time invariant, not BIBO stable, invertible. $T_2(x[n]) = x[n] + 0.5 x[n+1]$ is non-causal, linear, time invariant, BIBO stable, invertible.
  \item \textbf{Exercise 2:} Sketch $x_1[n] = 0.5^n u[n]$ and its shifts/reflections — preparation for visualizing convolution in the next lecture.
\end{enumerate}

\end{document}
